\documentclass[12pt]{article}
\usepackage[utf8]{inputenc}
\usepackage[T1]{fontenc}
\usepackage[margin=1in]{geometry}
\usepackage{setspace}
\usepackage{xcolor}
\usepackage[colorlinks=true, urlcolor=blue!60!black, linkcolor=blue!60!black]{hyperref}
\setlength{\parindent}{2em}
\setlength{\parskip}{1.25em}
\renewcommand{\baselinestretch}{1.0}
\title{\textit{Poindexter v. Greenhow}, 114 U.S. 270 (1885)}
\author{Case Brief}
\date{}
\begin{document}
\maketitle

\section{Case Metadata}
\begin{itemize}
  \item \textbf{Case Name:} \textit{Poindexter v. Greenhow} (lead decision of the \textit{Virginia Coupon Cases})
  \item \textbf{Citation:} \href{https://supreme.justia.com/cases/federal/us/114/270/}{114 U.S. 270} (\href{https://www.law.cornell.edu/supremecourt/text/114/270}{5 S. Ct. 903; 29 L. Ed. 185}) (1885)
  \item \textbf{Argued / Decided:} Argued March 20, 23--25, 1885; Decided April 20, 1885
  \item \textbf{Court and Posture:} Supreme Court of the United States; in error to the Hustings Court of the City of Richmond, Virginia
  \item \textbf{Briefing Book Placement:} Chapter 4: Federalism and State Sovereignty
  \item \textbf{Vote:} 5--4
  \item \textbf{Majority Author:} Justice Stanley Matthews (joined by Justices Field, Harlan, Woods, and Blatchford)
  \item \textbf{Dissenting Author:} Justice Joseph P. Bradley (joined by Chief Justice Waite and Justices Miller and Gray; published at \href{https://supreme.justia.com/cases/federal/us/114/325/}{114 U.S. 330} following \textit{Marye v. Parsons})
  \item \textbf{Constitutional Provisions:} U.S. Const. art. I, \S~10, cl. 1 (Bills of Credit Clause and Contracts Clause); U.S. Const. amend. XI (State Sovereign Immunity)
\end{itemize}

\section{Facts and Procedural History}
Facing crushing post-Civil War debt, Virginia enacted the Funding Act of March 30, 1871, authorizing creditors to exchange pre-war state obligations for new long-term bonds carrying detached semiannual interest coupons payable to bearer. To make the bonds marketable, the statute declared that the coupons would ``be receivable at and after maturity for all taxes, debts, dues, and demands due the state'' (\href{https://supreme.justia.com/cases/federal/us/114/270/}{114 U.S. at 270}).

When the populist Readjuster Party gained control of Virginia's General Assembly, it sought to curb coupon redemptions that were draining cash tax revenues. By the Act of January 26, 1882---later reinforced by the amendatory Act of March 13, 1884---Virginia forbade tax collectors from accepting coupons in payment of taxes, mandated payment solely in gold, silver, United States treasury notes, or national bank currency, and abolished common-law remedies (including detinue and trespass) against collectors. Taxpayers who tendered coupons were relegated to paying cash under protest and suing for verification and refund.

Poindexter, a Richmond taxpayer owing \$12.45 in state taxes, tendered matured 1871 coupons along with lawful currency for the fractional balance to Samuel C. Greenhow, Treasurer of the City of Richmond. Following the 1882 Act, Greenhow refused the coupons and, invoking Virginia's April 1, 1879 tax-distraint statute, seized Poindexter's office desk to satisfy the asserted delinquency (\href{https://www.law.cornell.edu/supremecourt/text/114/270}{114 U.S. at 273--74}). On April 26, 1883, Poindexter brought an action of detinue against Greenhow personally to recover the desk. A Richmond police justice dismissed the action for want of jurisdiction. On appeal, the Hustings Court of the City of Richmond---the highest state court with jurisdiction given the amount in controversy---held that jurisdiction existed, but entered judgment for Greenhow on the merits, ruling that the 1882 Act validly barred coupon tenders and confined Poindexter to its pay-under-protest procedure. Poindexter obtained a writ of error from the United States Supreme Court, which heard his case alongside seven companion \textit{Virginia Coupon Cases}.

\section{Issue(s) Presented}
\begin{enumerate}
  \item \textbf{Bills of Credit Clause:} Are detached, bearer-payable interest coupons issued under Virginia's 1871 Funding Act and made receivable for state taxes unconstitutional ``bills of credit'' prohibited by Article I, \S~10, cl. 1?
  \item \textbf{Contracts Clause:} Do Virginia's 1882 and 1884 statutes forbidding tax collectors to accept matured 1871 coupons and abolishing common-law tort remedies against collectors unconstitutionally impair the obligation of contracts under Article I, \S~10, cl. 1?
  \item \textbf{Eleventh Amendment:} Is a taxpayer's common-law detinue action against a state tax collector personally, seeking return of property seized after a lawful tender of coupons, barred by the Eleventh Amendment as a suit against the State of Virginia?
\end{enumerate}

\section{Holding}
\begin{enumerate}
  \item \textbf{No.} Tax-receivable state bond coupons are \emph{not} unconstitutional ``bills of credit'' because they become receivable only at maturity and are designed to assure redemption of state debt, not to circulate as money (\href{https://supreme.justia.com/cases/federal/us/114/270/}{114 U.S. at 283--85}).
  \item \textbf{Yes.} Virginia's 1871 Funding Act created a binding contract between the State and every coupon holder; the 1882 and 1884 Acts impaired that obligation in violation of the Contracts Clause and were void in their entirety, including their repeal of common-law remedies (\href{https://supreme.justia.com/cases/federal/us/114/270/}{114 U.S. at 276--82, 304--06}).
  \item \textbf{No.} Because a lawful tender of coupons is legally equivalent to payment, Greenhow lacked valid statutory authority to distrain Poindexter's desk; the detinue action therefore proceeded against Greenhow personally as a wrongdoer rather than against the State, falling outside the Eleventh Amendment (\href{https://supreme.justia.com/cases/federal/us/114/270/}{114 U.S. at 285--99}). Reversed and remanded.
\end{enumerate}

\section{Reasoning}
Writing for a five-Justice majority, Justice Matthews dismantled Virginia's constitutional and jurisdictional defenses in three steps.

First, Matthews rejected Virginia's attempt to void its own 1871 bonds under the Bills of Credit Clause. Drawing on \href{https://supreme.justia.com/cases/federal/us/29/410/}{\textit{Craig v. Missouri}}, 29 U.S. (4 Pet.) 410 (1830), and \href{https://supreme.justia.com/cases/federal/us/36/257/}{\textit{Briscoe v. Bank of the Commonwealth of Kentucky}}, 36 U.S. (11 Pet.) 257 (1837), the Court emphasized that a ``bill of credit'' requires paper issued by a state on its faith and credit \emph{and designed to circulate as money} in ordinary commercial transactions. Virginia's coupons lacked that monetary purpose: they were detached interest obligations tied to long-term bonds, became receivable for taxes only upon maturity, and were expected to be retired when due rather than passed from hand to hand as currency.

Second, on the Contracts Clause, the Court held that the 1871 Funding Act's tax-receivability pledge was ``of the essence of the right'' (\href{https://supreme.justia.com/cases/federal/us/114/270/}{114 U.S. at 271}). Because a lawful tender of taxes equals actual payment, Poindexter's tender extinguished Greenhow's authority to distrain. Virginia's substitute remedy---paying cash under protest and suing for a refund---did not preserve the contract; it destroyed the very right to pay in coupons. Nor could Virginia sever the provisions stripping taxpayers of detinue and trespass actions from the unconstitutional ban on coupon acceptance, because the entire statutory scheme was engineered to defeat the contract (\href{https://supreme.justia.com/cases/federal/us/114/270/}{114 U.S. at 304--06}).

Third, Matthews confronted the Eleventh Amendment. While acknowledging that sovereign immunity cannot be evaded by formalistic pleading, the Court distinguished suits seeking affirmative performance from the state treasury (barred under \textit{Louisiana v. Jumel}, 107 U.S. 711 (1883)) from tort actions against an officer who commits an affirmative wrong against a citizen's property under an unconstitutional statute. Invoking a sharp conceptual distinction between the \emph{State} (which can act only within constitutional bounds) and its \emph{government} (which can exceed those bounds), Matthews reasoned that an unconstitutional statute is ``as if it were not and never had been'' (\href{https://caselaw.findlaw.com/court/us-supreme-court/114/270.html}{114 U.S. at 292}). Stripped of his void statutory justification, Greenhow stood before the court as an ordinary trespasser answerable in detinue.

\section{Concurring and Dissenting Opinions}
No separate concurring opinions were filed. Justice Bradley dissented in an opinion joined by Chief Justice Waite and Justices Miller and Gray (reported at \href{https://supreme.justia.com/cases/federal/us/114/325/}{114 U.S. 330} after \textit{Marye v. Parsons}). Bradley did not defend Virginia's fiscal repudiation on moral grounds, nor did he dispute that the 1871 coupons formed a contract. Instead, he attacked the majority's Eleventh Amendment distinction as an unworkable legal fiction. In Bradley's view, a state's sovereign immunity bars all unconsented suits designed to coerce performance of state contracts, regardless of whether the plaintiff names the State or its tax collector. Allowing detinue or trespass suits against revenue officers who follow state law effectively strips the State of its prerogative to manage public debt and revenue, turning a forbidden suit for specific performance into a nominally personal tort action. Bradley also emphasized that Virginia's sovereign immunity applied equally against its own citizens, foreshadowing his landmark opinion five years later in \textit{Hans v. Louisiana}, 134 U.S. 1 (1890).

\section{Key Quotations}
\begin{quote}
``They do not become receivable in payment of taxes till they are due, and the design, we are bound to presume, was that they would be paid at maturity. This necessarily excludes the idea that they were intended for circulation at all.''\\
--- Justice Matthews, on the Bills of Credit Clause (\href{https://supreme.justia.com/cases/federal/us/114/270/}{114 U.S. at 284})
\end{quote}

\begin{quote}
``It is well settled by many decisions of this Court that, for the purpose of affecting proceedings to enforce the payment of taxes, a lawful tender of payment is equivalent to actual payment, either being sufficient to deprive the collecting officer of all authority for further action, and making every subsequent step illegal and void.''\\
--- Justice Matthews, on tender and official authority (\href{https://supreme.justia.com/cases/federal/us/114/270/}{114 U.S. at 282})
\end{quote}

\begin{quote}
``The defendant in error is not her officer, her agent, or her representative, in the matter complained of; for he has acted not only without her authority, but contrary to her express commands.''\\
--- Justice Matthews, articulating the officer-suit stripping fiction (\href{https://caselaw.findlaw.com/court/us-supreme-court/114/270.html}{114 U.S. at 292--93})
\end{quote}

\begin{quote}
``Unless the state chooses to allow itself to be sued, it cannot be sued. It has this prerogative if no other. ... In our judgment, none of these suits can be maintained, for the reason that they are, in substance and effect, suits against the State of Virginia.''\\
--- Justice Bradley, dissenting (\href{https://supreme.justia.com/cases/federal/us/114/325/}{114 U.S. at 335--36})
\end{quote}

\section{Doctrinal Significance and Modern Legacy (Through 2026)}
\textit{Poindexter} sits at the crossroads of three enduring federalism doctrines. First, its Bills of Credit holding---later reinforced in \href{https://www.law.cornell.edu/supremecourt/text/177/66}{\textit{Houston \& Texas Central Railroad Co. v. Texas}}, 177 U.S. 66 (1900)---cemented a narrow, intent-based definition of Article I, \S~10's monetary bar, ensuring states could freely issue municipal and state bonds, tax-anticipation warrants, and coupons without triggering federal invalidation.

Second, \textit{Poindexter} is a foundational pillar of the constitutional officer-suit exception. Its ``stripping'' theory---that an officer enforcing an unconstitutional state statute loses sovereign immunity---was refined in \href{https://www.law.cornell.edu/supremecourt/text/123/443}{\textit{In re Ayers}}, 123 U.S. 443 (1887) (barring suits that merely seek to compel state contract performance), and culminated in \href{https://supreme.justia.com/cases/federal/us/209/123/}{\textit{Ex parte Young}}, 209 U.S. 123 (1908), which directly cited \textit{Poindexter} to authorize prospective federal injunctions against state officials enforcing unconstitutional laws. Through 2026, the Supreme Court continues to police the boundary \textit{Poindexter} helped draw: officer suits survive for prospective relief against ongoing federal-law violations (\href{https://supreme.justia.com/cases/federal/us/458/670/}{\textit{Florida Dept. of State v. Treasure Salvors, Inc.}}, 458 U.S. 670 (1982)), but cannot enforce state law against state officers (\href{https://supreme.justia.com/cases/federal/us/465/89/}{\textit{Pennhurst State School \& Hospital v. Halderman}}, 465 U.S. 89 (1984)), invade special state sovereignty interests (\href{https://supreme.justia.com/cases/federal/us/521/261/}{\textit{Idaho v. Coeur d'Alene Tribe of Idaho}}, 521 U.S. 261 (1997)), override expansive sovereign immunity under \href{https://supreme.justia.com/cases/federal/us/517/44/}{\textit{Seminole Tribe of Florida v. Florida}}, 517 U.S. 44 (1996), and \href{https://supreme.justia.com/cases/federal/us/527/706/}{\textit{Alden v. Maine}}, 527 U.S. 706 (1999), or enjoin state judges and clerks where a state delegates enforcement exclusively to private litigants (\href{https://supreme.justia.com/cases/federal/us/595/21-463/}{\textit{Whole Woman's Health v. Jackson}}, 595 U.S. 30 (2021)).

Third, while modern Contracts Clause review has softened since the nineteenth century into a two-step ``substantial impairment'' and reasonableness inquiry---reaffirmed in \href{https://supreme.justia.com/cases/federal/us/584/16-1432/}{\textit{Sveen v. Melin}}, 584 U.S. 811 (2018)---courts still apply heightened scrutiny when a state abrogates its own financial obligations, preserving \textit{Poindexter}'s core anti-repudiation principle.

\section{Pedagogical Critiques and Exam Hypotheticals}
\textbf{Pedagogical Critique:} Students should recognize the central tension in \textit{Poindexter}: the majority insists that Greenhow is \emph{not} acting for Virginia under the Eleventh Amendment, yet his conduct violates the Contracts Clause---which restricts only ``State'' action---precisely because he enforces Virginia's 1882 Act. That ``state-actor-for-the-merits, private-wrongdoer-for-immunity'' paradox passed directly into \textit{Ex parte Young} and remains the defining analytical fiction of federal courts doctrine.

\textbf{``Opposite Side'' Contrast (\href{https://supreme.justia.com/cases/federal/us/29/410/}{\textit{Craig v. Missouri}}, 29 U.S. (4 Pet.) 410 (1830)):} Pair \textit{Poindexter} with \textit{Craig} to mark the exact constitutional boundary of the Bills of Credit Clause. In \textit{Craig}, Missouri issued small-denomination ``loan office certificates'' (\$0.50 to \$10) receivable for taxes and public salaries, backed by state faith, and explicitly structured to circulate as a substitute for specie; Chief Justice Marshall struck them down as forbidden bills of credit. In \textit{Poindexter}, Virginia's tax-receivable bond coupons stood on the opposite side of the line: tied to funded debt and receivable only at maturity, they functioned as debt-redemption assurances rather than circulating currency.

\textbf{Exam Hypothetical:} Facing a severe budget shortfall, State X issues \$10 million in bearer ``Tax Credit Notes'' in \$5, \$10, and \$20 denominations, payable immediately upon issuance, transferable by delivery, and receivable for all state taxes and fees. Two years later, State X repeals tax-receivability and directs county treasurers to seize property of taxpayers who tender the notes. When a taxpayer sues the county treasurer in federal court for return of a seized truck, how should the court rule? \textit{Analysis:} Under \textit{Craig} and \textit{Poindexter}, the threshold issue is whether the notes are void \textit{ab initio} as unconstitutional bills of credit. Unlike Virginia's matured bond coupons in \textit{Poindexter}, State X's small-denomination, immediately receivable bearer notes closely resemble Missouri's circulating loan certificates in \textit{Craig}. If void as bills of credit, no valid contract existed for the Contracts Clause to protect; if valid debt instruments, however, \textit{Poindexter} permits the officer suit against the treasurer personally unless framed as an affirmative raid on the state treasury.

\end{document}
